% Outline: Conclusion (exactly one paragraph)
\section{Conclusion}

Agent traffic carries a signal the serving layer currently discards: the
tool-schema text attached to every tool-bearing request. Read as text, the
request improves output-length ranking by $+0.19$ Kendall's $\tau$ over the
tuned scalar family we benchmark, and the schema itself contributes $+0.044$
of that; ranking quality changes little on the two adequately powered strata
of tool sets unseen in training (equivalence untested), where identity
features have nothing to look up. Under the gateway's documented 15\,ms decision timeout the predictor
silently loses 34.6\% of its decisions; the two CPU rescues we
pre-registered do not repair that, while a 75\,ms timeout eliminated logged
fail-open in these runs, permitting up to that much critical-path decision
latency on each scored request. Whether acting on the signal
pays depends on how contended the engine is, and that turned out to be a
configuration we set rather than a property of the traffic. With the
running-batch cap at its default the queue stayed empty at the 90th
percentile and no ordering
policy separated from another; capping the batch at sixteen put most requests
behind another, and there the learned ranker beat arrival order by 4.2\%
while ordering by prompt length ran 8.5\% slower than leaving the queue
alone. The reliability gate follows the same regime split: free when nothing
waits, and 9.5\% expensive once something does, because the requests it
declines to vouch for hold their arrival slots and block. Two levers paid in
the uncapped configuration and neither came from prediction: exchanging the stock scheduler
for the policy implementation cut mean time-to-last-token by 44\% with
arrival order held fixed, and prefix caching over the resent schemas cut a
further 17--23\%. The schema signal is real, the cheap scalar substitute for
it is worse than doing nothing, and the operator who wants either had better
check what the batch cap leaves for a scheduler to decide.
